\documentclass[11pt,a4paper]{article}

\usepackage[utf8]{inputenc}
\usepackage[T1]{fontenc}
\usepackage{amsmath,amssymb}
\usepackage{booktabs}
\usepackage{geometry}
\usepackage{algorithm}
\usepackage{algpseudocode}
\usepackage[numbers,sort&compress]{natbib}
\usepackage{hyperref}
\usepackage{xcolor}

\geometry{margin=2.6cm}
\hypersetup{colorlinks=true,linkcolor=blue!50!black,citecolor=blue!50!black,
            urlcolor=blue!50!black}

\title{Koiter asymptotic analysis about a non-linear pre-buckling state\\
\large Implementation notes for \texttt{NLprebuck=True} in
\texttt{bfsccylinder\_models}}
\author{Saullo G. P. Castro}
\date{}

\begin{document}
\maketitle

\begin{abstract}
This note documents the implementation of the single-mode Koiter method with a
non-linear pre-buckling state in
\texttt{koiter\_cylinder\_newton\_raphson.py} (Donnell-Mushtari-Vlasov
kinematics) and \texttt{koiter\_cylinder\_newton\_raphson\_sanders.py}
(Sanders-Koiter kinematics), and of the same algorithm in the
variable-stiffness models \texttt{koiter\_cylinder\_CTS.py} and
\texttt{koiter\_cylinder\_CTS\_sanders.py} (Section~\ref{sec:cts}).
It covers the tangent stiffness matrix and the
consistency it must satisfy, the reduction of the pre-buckling problem to the
axisymmetric subspace, the incremental-iterative scheme that places the
expansion point in the neighbourhood of the bifurcation point, the removal of
the arbitrariness left by the degenerate buckling modes, and the
stiffness-weighted orthogonality condition imposed on the second-order
fields, and the mesh resolution the continuous-tow-shearing cylinders
require.
The notation follows the finite element chapter of \citet{castro2023}.  Equation
numbers of the form ``Sun~Eq.~(nn)'' refer to \citet{sun2020}.
\end{abstract}

\tableofcontents

%==============================================================================
\section{Scope and notation}
%==============================================================================

The routine \texttt{fkoiter\_cyl\_SS3} performs, for a circular cylindrical
shell under axial compression with SS-3 boundary conditions,
\begin{enumerate}\itemsep2pt
  \item a pre-buckling analysis,
  \item a buckling eigenvalue analysis about the pre-buckling state,
  \item the Koiter asymptotic expansion about the bifurcation point, returning
        the post-buckling coefficients $a_{ijk}$ and $b_{ijkl}$.
\end{enumerate}

With \texttt{NLprebuck=False} the pre-buckling state is the linear elastic
solution and the pre-buckling path is exactly linear in the load parameter.
With \texttt{NLprebuck=True} the pre-buckling state is obtained from a
geometrically non-linear analysis, and the whole expansion must then be
carried out about a state \emph{in the neighbourhood of the bifurcation
point}.  This document describes the latter case.

The shell coordinates are the axial $x$ and the circumferential $y$, the
radius is $r$ and the thickness $h$.  Following \citet{dossantos2022}, $y$
is the circumferential coordinate measured as an \emph{arc length},
$y = r\vartheta$ with $\vartheta$ the polar angle, so that $x$ and $y$
have the same units and the membrane strains below carry no explicit
factor of $r$.  The symbol $\theta$ is reserved throughout for the tow
steering angle of Section~\ref{sec:cts}.  The circumferential
discretisation is $n_y$ elements around the circumference, the axial one
$n_x$ stations along the generator.  The degrees-of-freedom are collected in
$\boldsymbol{\bar{u}}$, the fundamental (pre-buckling) state is
$\boldsymbol{\hat{u}}$, and the buckling modes are $\delta\boldsymbol{\bar{u}}_i$.
The external load is proportional, $\boldsymbol{F}_{ext} = \lambda
\boldsymbol{F}_{ref}$, with $\boldsymbol{F}_{ref}$ the reference load vector
built from $N_{xx}^{\text{unit}}$.  Two load levels are distinguished:
$\lambda_b$, at which the pre-buckling state is evaluated, and $\lambda_c$,
the bifurcation load.

%==============================================================================
\section{Kinematics}
%==============================================================================

The membrane strains are written, for both kinematic models, as
\begin{equation}\label{eq:strain}
	\boldsymbol{\varepsilon}
	= \boldsymbol{B}_m\boldsymbol{\bar{u}}
	+ \frac{1}{2}
	\begin{Bmatrix} g_x^2 \\ g_y^2 \\ 2\,g_x g_y \end{Bmatrix},
	\qquad
	\boldsymbol{\kappa} = \boldsymbol{B}_b\boldsymbol{\bar{u}}
\end{equation}

\noindent where $g_x$ and $g_y$ are the two gradient measures that carry
the geometric non-linearity:
\begin{equation}\label{eq:gradients}
	\text{Donnell-Mushtari-Vlasov:}\quad
	g_x = w_{,x},\quad g_y = w_{,y},
	\qquad
	\text{Sanders-Koiter:}\quad
	g_x = w_{,x},\quad g_y = w_{,y} - \frac{v}{r}
\end{equation}

Every non-linear term of the Sanders-Koiter membrane strains is recovered from
the Donnell-Mushtari-Vlasov one by the substitution
$(w_{,x},w_{,y})\rightarrow(g_x,g_y)$, which is why the two modules
differ only in the definition of Eq.~\ref{eq:gradients} and in the
corresponding interpolation matrices $\boldsymbol{G}_x$ and
$\boldsymbol{G}_y$, with $g_x = \boldsymbol{G}_x\boldsymbol{\bar{u}}$ and
$g_y = \boldsymbol{G}_y\boldsymbol{\bar{u}}$.  The additional term
$-v/r$ of the Sanders-Koiter operator $\boldsymbol{G}_y$ is the same one
highlighted in the geometric stiffness matrix of the linear buckling chapter.
The constitutive relation is the usual laminated-shell one, collectively
denoted $\boldsymbol{C}$, relating $\{\boldsymbol{N},\boldsymbol{M}\}$ to
$\{\boldsymbol{\varepsilon},\boldsymbol{\kappa}\}$ through the $ABD$ matrix.

\subsection{State and rates are independent fields}
\label{sec:staterates}

The pre-buckling state is a curve $\boldsymbol{\hat{u}}(\lambda)$.  The
asymptotic expansion needs, besides the state itself, its first and second
derivatives with respect to $\lambda$, written $\boldsymbol{\hat{u}}_{,\lambda}$ and
$\boldsymbol{\hat{u}}_{,\lambda\lambda}$.  Differentiating Eq.~\ref{eq:strain} along the
path, and abbreviating $\hat{g}_\alpha =
\boldsymbol{G}_\alpha\boldsymbol{\hat{u}}$, $g_{\alpha,\lambda} =
\boldsymbol{G}_\alpha\boldsymbol{\hat{u}}_{,\lambda}$, $g_{\alpha,\lambda\lambda} =
\boldsymbol{G}_\alpha\boldsymbol{\hat{u}}_{,\lambda\lambda}$, with
$\alpha \in \{x,y\}$:
\begin{align}
	\boldsymbol{\varepsilon}_{,\lambda}
	&= \boldsymbol{B}_m\boldsymbol{\hat{u}}_{,\lambda}
	 + \begin{Bmatrix}
	     \hat{g}_x g_{x,\lambda} \\ \hat{g}_y g_{y,\lambda} \\
	     \hat{g}_x g_{y,\lambda} + \hat{g}_y g_{x,\lambda}
	   \end{Bmatrix},
	& \boldsymbol{\kappa}_{,\lambda} &= \boldsymbol{B}_b\boldsymbol{\hat{u}}_{,\lambda}
	\label{eq:epsdot}\\[4pt]
	\boldsymbol{\varepsilon}_{,\lambda\lambda}
	&= \boldsymbol{B}_m\boldsymbol{\hat{u}}_{,\lambda\lambda}
	 + \begin{Bmatrix}
	     g_{x,\lambda}^2 + \hat{g}_x g_{x,\lambda\lambda} \\
	     g_{y,\lambda}^2 + \hat{g}_y g_{y,\lambda\lambda} \\
	     2 g_{x,\lambda} g_{y,\lambda} + \hat{g}_x g_{y,\lambda\lambda}
	       + \hat{g}_y g_{x,\lambda\lambda}
	   \end{Bmatrix},
	& \boldsymbol{\kappa}_{,\lambda\lambda} &= \boldsymbol{B}_b\boldsymbol{\hat{u}}_{,\lambda\lambda}
	\label{eq:epsddot}
\end{align}

\noindent and the derivatives with respect to the degrees-of-freedom:
\begin{equation}\label{eq:epsderivs}
	\frac{\partial\boldsymbol{\varepsilon}}{\partial\boldsymbol{\bar{u}}}
	= \boldsymbol{B}_m + \begin{Bmatrix}
	    \hat{g}_x\boldsymbol{G}_x \\ \hat{g}_y\boldsymbol{G}_y \\
	    \hat{g}_x\boldsymbol{G}_y
	      + \hat{g}_y\boldsymbol{G}_x \end{Bmatrix},
	\qquad
	\frac{\partial^2\boldsymbol{\varepsilon}}
	     {\partial\lambda\,\partial\boldsymbol{\bar{u}}}
	= \begin{Bmatrix}
	    g_{x,\lambda}\boldsymbol{G}_x \\ g_{y,\lambda}\boldsymbol{G}_y \\
	    g_{x,\lambda}\boldsymbol{G}_y
	      + g_{y,\lambda}\boldsymbol{G}_x \end{Bmatrix}
\end{equation}

If one assumes the pre-buckling path to be linear in the load parameter,
$\boldsymbol{\hat{u}}(\lambda) = \lambda\,\boldsymbol{\hat{u}}(1)$, then
$\boldsymbol{\hat{u}}_{,\lambda} = \boldsymbol{\hat{u}}(1)$,
$\boldsymbol{\hat{u}}_{,\lambda\lambda} = \boldsymbol{0}$ and $\hat{g}_\alpha =
\lambda\,g_{\alpha,\lambda}$, so Eqs.~\ref{eq:epsdot} to \ref{eq:epsderivs}
collapse onto expressions containing $\lambda$ explicitly.  That assumption is
precisely what removes the non-linear pre-buckling behaviour from the
expansion, and it is not made here: state, first rate and second rate are
carried as three independent nodal fields.

%==============================================================================
\section{The tangent stiffness matrix}
\label{sec:tangent}
%==============================================================================

\subsection{Definition and decomposition}

The stiffness matrix used throughout is the derivative of the internal force
vector with respect to the degrees-of-freedom.  Differentiating the residual
$\boldsymbol{R} = \boldsymbol{F}_{int} - \boldsymbol{F}_{ext}$, and with no
follower forces, so that $\partial\boldsymbol{F}_{ext}/
\partial\boldsymbol{\bar{u}} = \boldsymbol{0}$:
\begin{equation}\label{eq:tangent-definition}
	\boldsymbol{R}' = \frac{\partial \boldsymbol{R}}{\partial \boldsymbol{\bar{u}}}
	= \frac{\partial \boldsymbol{F}_{int}}{\partial \boldsymbol{\bar{u}}}
	= \boldsymbol{K}_T
\end{equation}

Differentiating the internal force vector produces one contribution from the
change of the stresses and another from the change of the differentiation
operator itself:
\begin{equation}\label{eq:tangent-decomposition}
	\boldsymbol{K}_T
	= \underbrace{\int_{\Omega} (\boldsymbol{B}_L + \boldsymbol{B}_{NL})^\top
	  \boldsymbol{C} (\boldsymbol{B}_L + \boldsymbol{B}_{NL})
	  \, d\Omega}_{\boldsymbol{K}_C}
	+ \underbrace{\int_{\Omega}
	  \frac{\partial (\boldsymbol{B}_L + \boldsymbol{B}_{NL})^\top}
	       {\partial \boldsymbol{\bar{u}}}
	  \boldsymbol{\sigma} \, d\Omega}_{\boldsymbol{K}_G}
	= \boldsymbol{K}_C + \boldsymbol{K}_G
\end{equation}

For the elements used here the constitutive contribution is split further into
the constant part and the part that depends on the displacement state:
\begin{equation}\label{eq:KC-split}
	\boldsymbol{K}_C(\boldsymbol{\bar{u}})
	= \boldsymbol{K}_0 + \boldsymbol{K}_C^{NL}(\boldsymbol{\bar{u}}),
	\qquad
	\boldsymbol{K}_0 = \int_{\Omega} \boldsymbol{B}_L^\top \boldsymbol{C}
	                   \boldsymbol{B}_L \, d\Omega
\end{equation}

\noindent which matches the three assembly routines of the element library,
\texttt{update\_KC0}, \texttt{update\_KCNL} and \texttt{update\_KG}, so that
$\boldsymbol{K}_T = \boldsymbol{K}_0 + \boldsymbol{K}_C^{NL} +
\boldsymbol{K}_G$.  The internal force vector is assembled separately, by
\texttt{update\_fint}.

\subsection{Verifying consistency}
\label{sec:tangent-check}

Nothing in the assembly guarantees that the three matrices add up to the
derivative of the separately assembled $\boldsymbol{F}_{int}$; that has to be
verified.  The test is a directional Taylor expansion.  For an arbitrary
direction $\boldsymbol{d}$ and a step $\epsilon$:
\begin{equation}\label{eq:taylor-test}
	e(\epsilon) = \frac{\left\lVert \boldsymbol{F}_{int}
	  (\boldsymbol{\bar{u}} + \epsilon\boldsymbol{d})
	  - \boldsymbol{F}_{int}(\boldsymbol{\bar{u}})
	  - \epsilon\,\boldsymbol{K}_T\boldsymbol{d} \right\rVert}
	  {\left\lVert \epsilon\,\boldsymbol{K}_T\boldsymbol{d} \right\rVert}
\end{equation}

The numerator is $O(\epsilon^2)$ when $\boldsymbol{K}_T$ is the derivative, so
$e(\epsilon)$ must fall proportionally to $\epsilon$, dividing by ten at every
decade, until round-off takes over.  A plateau means the matrix is not the
derivative.  Measured on a single element with a displacement state of the
order of the thickness:

\begin{center}
\begin{tabular}{lcc}
\toprule
$\epsilon$ & before \texttt{bfsccylinder} 0.6.0 & from 0.6.0 onwards \\
\midrule
$10^{-2}$ & $6.186\times10^{-2}$ & $9.466\times10^{-4}$ \\
$10^{-3}$ & $6.157\times10^{-2}$ & $9.470\times10^{-5}$ \\
$10^{-4}$ & $6.154\times10^{-2}$ & $9.470\times10^{-6}$ \\
$10^{-5}$ & $6.154\times10^{-2}$ & $9.470\times10^{-7}$ \\
$10^{-6}$ & $6.154\times10^{-2}$ & $9.450\times10^{-8}$ \\
\bottomrule
\end{tabular}
\end{center}

A second test identifies where a deficiency sits.  Taking the true directional
derivative $\boldsymbol{J}\boldsymbol{d}$ from a central difference and fitting
the coefficients of the two state-dependent matrices by least squares,
\begin{equation}\label{eq:tangent-fit}
	\min_{c_1,c_2} \left\lVert \boldsymbol{K}_0\boldsymbol{d}
	  + c_1 \boldsymbol{K}_C^{NL}\boldsymbol{d}
	  + c_2 \boldsymbol{K}_G\boldsymbol{d}
	  - \boldsymbol{J}\boldsymbol{d} \right\rVert
\end{equation}

\noindent returns $c_1 = c_2 = 1.0000$ with a residual of $7.6\times10^{-11}$
from 0.6.0 onwards.  Before 0.6.0 it returned $c_1 = 1.5122$ and $c_2 =
1.2876$ with a residual of $3.8\times10^{-2}$, against $6.0\times10^{-2}$ for
the assembled combination: no recombination of the three matrices reproduced
the derivative, which identified the defect as a missing term rather than a
mis-scaled one.  Both kinematics are covered, the Sanders-Koiter element
reaching $1.2\times10^{-10}$.

\subsection{Why the consistency matters twice}

An inconsistent tangent stiffness matrix has two distinct consequences here,
and only the first is the familiar one.

In the Newton-Raphson solution of Section~\ref{sec:axisymmetric} the tangent
controls the rate of convergence only, the converged state being fixed
exclusively by the residual, that is by the accuracy with which
$\boldsymbol{F}_{int}$ and $\boldsymbol{F}_{ext}$ are computed.  An
inconsistent tangent therefore degrades the iteration from quadratic to
linear, with a contraction factor that grows as the tangent is approached, but
it does not by itself move the equilibrium state.  On the coarse meshes of
Section~\ref{sec:results} the contraction was about $0.3$ per iteration, which
still converges; on a mesh with $n_y = 60$ it reached $0.95$, exhausting
the iteration budget at every load step and leaving the load stepping stalled
near $\lambda_b/\lambda_c = 0.94$.

The second consequence has no such protection.  The matrix
$\boldsymbol{K}_C = \boldsymbol{K}_0 + \boldsymbol{K}_C^{NL}$ does not only
drive the iteration: it enters the eigenvalue problem of
Eq.~\ref{eq:eigenproblem} and the expansion operator of
Eq.~\ref{eq:phi2} directly, where it is not a preconditioner but part of the
statement of the problem.  An error in $\boldsymbol{K}_C^{NL}$ is therefore
inherited by $\lambda_c$, by the buckling mode and by every post-buckling
coefficient.  Correcting it moved the buckling loads of
Section~\ref{sec:results} by $9$ to $11$ per cent, onto the reference
solutions, and reversed the sign of $b_{1111}$.

%==============================================================================
\section{The axisymmetric pre-buckling solver}
\label{sec:axisymmetric}
%==============================================================================

\subsection{The fundamental path is axisymmetric}

The fundamental path of a perfect cylinder under uniform axial compression is
axisymmetric, and every buckling mode that makes $\boldsymbol{K}_T$ singular
is not.  Solving the pre-buckling problem in the axisymmetric subspace
therefore keeps the iteration clear of the entire near-critical cluster, and
not merely of the few modes that happen to have been computed.  This is also
what the reference solutions do, AXBIF and ANILISA both solving an
axisymmetric pre-buckling problem.

\subsection{Cyclic symmetry of the mesh}

The mesh has $n_y$ elements around the circumference, and the element
degrees-of-freedom are expressed in the local shell frame, which turns with
$y$.  Shifting the circumferential node index cyclically by one element
is therefore a symmetry of every operator assembled on the mesh: if
$\boldsymbol{P}$ denotes that shift, then $\boldsymbol{P}^\top
\boldsymbol{K}\boldsymbol{P} = \boldsymbol{K}$ for
$\boldsymbol{K} \in \{\boldsymbol{K}_0, \boldsymbol{K}_C^{NL},
\boldsymbol{K}_G\}$ whenever the state they are evaluated at is itself
invariant.

Two consequences are used.  The first is that the average over the orbit of
$\boldsymbol{P}$ is the orthogonal projector onto the invariant subspace, used
to remove the round-off asymmetry of the starting state.  The second, used in
Section~\ref{sec:canonical}, is that $\boldsymbol{P}$ maps a buckling mode onto
another buckling mode carrying the same multiplier.

\subsection{The reduced problem}

Let $\boldsymbol{B}_{axi}$ collect one column per axial station and per
degree-of-freedom, carrying one at that degree-of-freedom of every node of
the station,
\begin{equation}\label{eq:baxi}
	\left[\boldsymbol{B}_{axi}\right]_{\,(i\,n_y + j)\,n_{DOF} + d,\;\;
	                                    i\,n_{DOF} + e} = \delta_{de},
	\qquad
	\begin{array}{l}
	i = 0\ldots n_x-1 \\ j = 0\ldots n_y-1 \\ d,e = 0\ldots n_{DOF}-1
	\end{array}
\end{equation}
so that an axisymmetric field is $\boldsymbol{B}_{axi}\boldsymbol{a}$, the
columns are mutually orthogonal, and $\boldsymbol{B}_{axi}^\top
\boldsymbol{B}_{axi} = n_y\boldsymbol{I}$.

The correction is then obtained by a Bubnov-Galerkin projection
\citep{fletcher1984}: it is sought \emph{in} the subspace spanned by the
columns of $\boldsymbol{B}_{axi}$, and the residual of the tangent problem is
required to be orthogonal \emph{to that same subspace},
\begin{equation}\label{eq:galerkin}
	\delta\boldsymbol{\bar{u}} = \boldsymbol{B}_{axi}\boldsymbol{a}
	\quad\text{(trial)},
	\qquad
	\boldsymbol{B}_{axi}^\top\left(\boldsymbol{K}_T\,
	\boldsymbol{B}_{axi}\boldsymbol{a} + \boldsymbol{R}\right)
	= \boldsymbol{0}
	\quad\text{(test)}
\end{equation}
Trial and test spaces coincide, which is what makes the projection
Bubnov-Galerkin rather than Petrov-Galerkin, and which preserves the symmetry
of $\boldsymbol{K}_T$ in the reduced operator.  Written out, Eq.
\ref{eq:galerkin} is
\begin{equation}\label{eq:reduced-newton}
	\left( \boldsymbol{B}_{axi}^\top \boldsymbol{K}_T \boldsymbol{B}_{axi}
	\right) \boldsymbol{a}
	= -\,\boldsymbol{B}_{axi}^\top \boldsymbol{R},
	\qquad
	\delta\boldsymbol{\bar{u}} = \boldsymbol{B}_{axi}\boldsymbol{a}
\end{equation}

This is the reduced-basis technique for non-linear structural analysis of
\citet{noor1980}, used here with a basis known \emph{a priori} from the
symmetry of the problem instead of one accumulated from previous solutions.
The cyclic shift $\boldsymbol{P}$ commutes with $\boldsymbol{K}_T$, and the
column space of $\boldsymbol{B}_{axi}$ is exactly the subspace invariant
under it; building the basis from the symmetry group in this way is the
standard treatment of rotationally periodic structures \citep{thomas1979}.
What is specific to the present use is that only the zeroth circumferential
harmonic is retained, the remaining $n_y - 1$ harmonics being precisely where
the buckling modes live.

Three consequences follow.

\begin{enumerate}\itemsep3pt
  \item \textbf{Size.}  The reduced system has $n_x \times n_{DOF}$ unknowns
        against the full count, $250$ against $15000$ on the $n_y = 60$
        mesh, and is small enough to be solved as a dense system.
  \item \textbf{No deflation.}  The whole near-critical cluster lies outside
        the subspace by construction, so the buckling modes never have to be
        deflated out of the correction.
  \item \textbf{The reduced solution is exact.}  Because $\boldsymbol{K}_T$
        and $\boldsymbol{R}$ commute with $\boldsymbol{P}$, the residual
        evaluated at an axisymmetric state is itself axisymmetric and
        therefore has no component outside the subspace.  Hence
        $\boldsymbol{B}_{axi}^\top\boldsymbol{R} = \boldsymbol{0}$ implies
        $\boldsymbol{R} = \boldsymbol{0}$ on the free degrees-of-freedom, and
        the Galerkin condition of Eq.~\ref{eq:galerkin} is \emph{equivalent}
        to the full equilibrium equations rather than a projection of them.
        The converged state solves the full problem exactly.  This is what
        separates a symmetry-adapted basis from a general reduced basis, for
        which Eq.~\ref{eq:galerkin} leaves a residual outside the subspace
        that is never driven to zero.
\end{enumerate}

A reduced coordinate is treated as known as soon as one of the nodes it spans
is constrained.  This matters for the single-node constraint that suppresses
the axial rigid body translation, which is the one part of the boundary
conditions that is not rotationally symmetric: in the reduced problem it fixes
the whole axial station, which is what an axisymmetric field requires of it
anyway.  The same asymmetry is handled in Section~\ref{sec:canonical} by
subtracting an axial rigid body translation, which carries no strain and is
annihilated by both stiffness matrices.

\subsection{Conditioning of the reduced matrix}
\label{sec:scaling}

The reduced coordinates place nodal displacements next to nodal derivatives,
whose stiffness entries differ by many orders of magnitude, and the free part
of $\boldsymbol{K}_r = \boldsymbol{B}_{axi}^\top \boldsymbol{K}_T
\boldsymbol{B}_{axi}$ reaches a condition number of $4\times10^{12}$ on the
$n_y = 60$ mesh.  The sparse solver used for the full systems equilibrates
internally; the dense solver used on the reduced one does not, so the
symmetric (Jacobi) equilibration is applied explicitly.

Let $\boldsymbol{D} = \mathrm{diag}\left(|K_{r,kk}|^{-1/2}\right)$.  Then
$\boldsymbol{K}_r\boldsymbol{a} = \boldsymbol{r}$ is solved in the form
\begin{equation}\label{eq:scaling}
	\left(\boldsymbol{D}\boldsymbol{K}_r\boldsymbol{D}\right)
	\left(\boldsymbol{D}^{-1}\boldsymbol{a}\right)
	= \boldsymbol{D}\boldsymbol{r},
	\qquad\text{so}\qquad
	\boldsymbol{a} = \boldsymbol{D}\left[
	\left(\boldsymbol{D}\boldsymbol{K}_r\boldsymbol{D}\right)^{-1}
	\boldsymbol{D}\boldsymbol{r}\right]
\end{equation}
which is what the code implements: it forms
$\boldsymbol{D}\boldsymbol{K}_r\boldsymbol{D}$ and
$\boldsymbol{D}\boldsymbol{r}$, solves the scaled system, and multiplies the
result by $\boldsymbol{D}$ once more.  The scaled matrix has unit diagonal in
absolute value.  For a symmetric positive definite matrix, equilibrating the
diagonal in this way is within a factor of $n$ of the optimal symmetric
scaling \citep{vandersluis1969}, $n$ being the order of the matrix; the
reduced tangent is positive definite everywhere the load stepping actually
solves it, which is on the near side of the bifurcation point.  The measured
condition number falls from $4\times10^{12}$ to $3\times10^{2}$.

Two points are worth making explicit.  Eq.~\ref{eq:scaling} is an exact
rewriting of the same linear system, not a preconditioner and not an
approximation: it changes the floating-point path taken, never the solution
sought.  And on the meshes tried it changes nothing measurable --- on the
$n_y = 60$ case every iterate is unchanged to ten digits, the graded matrix
being solved accurately by partial pivoting alone.  It is kept as a guard
against worse-conditioned meshes, not as a correction to these.

%==============================================================================
\section{Locating the expansion point}
\label{sec:expansion-point}
%==============================================================================

\subsection{Why the load level matters}

Koiter's expansion is asymptotic about the bifurcation point.  Sun~\S2.2
requires the pre-buckling state to satisfy $\lambda_b/\lambda_c \approx
0.995$.  If instead the pre-buckling problem is solved once at the reference
load and the eigenvalue problem is posed about that state, then for the shells
considered here $\lambda_b/\lambda_c$ is of the order of $0.1$.  At that load
level the radial pre-buckling deflection is one to three per cent of a wall
thickness, the boundary layer produced by the edge restraint has not
developed, and $\lVert\boldsymbol{K}_C^{NL}(\boldsymbol{\hat{u}})\rVert /
\lVert\boldsymbol{K}_0\rVert \sim 2\times10^{-4}$; the eigenvalue analysis is
then a linear buckling analysis in all but name and the whole computation
returns the \emph{membrane} pre-buckling answer.

\subsection{Iterative eigenvalue algorithm}

The expansion point is located by the incremental-iterative scheme of
Sun~Eqs.~(44) to (46).  At the current load level $\lambda_b$ the pre-buckling
state $\boldsymbol{\hat{u}}$ is known and the eigenvalue problem
\begin{equation}\label{eq:eigenproblem}
	\left[ \boldsymbol{K}_C(\boldsymbol{\hat{u}})
	+ \mu\,\boldsymbol{K}_G(\boldsymbol{\hat{u}}) \right]
	\delta\boldsymbol{\bar{u}} = \boldsymbol{0}
\end{equation}

\noindent gives the multiplier $\mu$ \emph{of the current stress state}, so
that $\lambda_c = \lambda_b\,\mu$.  This is the eigenvalue problem of the
linear buckling chapter with $\boldsymbol{K}_0$ replaced by the constitutive
part of the tangent at the current state.

\subsection{Advancing the load level}
\label{sec:step}

Sun~Eq.~(44) advances the load level by a fixed fraction of the distance to
the bifurcation load:
\begin{equation}\label{eq:step}
	\lambda_b \leftarrow \lambda_b + \eta\,(\lambda_c - \lambda_b)
\end{equation}

\noindent with $\eta = 0.8$ on the first step and $\eta = 0.5$ afterwards.
The difficulty is that $\lambda_c$ is not fixed: it is re-evaluated at the new
state and it falls as $\lambda_b$ rises.  Writing $s = d\lambda_c/d\lambda_b$
and linearising, the new distance to the bifurcation point is
\begin{equation}\label{eq:step-overshoot}
	\lambda_b' - \lambda_c'
	= (\lambda_c - \lambda_b)\left[ \eta(1-s) - 1 \right]
\end{equation}

\noindent so the step lands \emph{past} the bifurcation point whenever
\begin{equation}\label{eq:eta-critical}
	\eta > \frac{1}{1-s}
\end{equation}

The sensitivity $s$ steepens as the bifurcation point is approached and as the
mesh is refined.  On the $n_y = 40$ meshes of Section~\ref{sec:results}
it reaches $-0.25$ and $-0.72$, for which Eq.~\ref{eq:eta-critical} permits
$\eta = 0.80$ and $\eta = 0.58$, both above the fixed $0.5$, and no overshoot
occurs.  On the $n_y = 60$ mesh it reaches $-1.09$, permitting only
$\eta = 0.478$, and the fixed $\eta = 0.5$ overshoots to
$\lambda_b/\lambda_c = 1.0013$.

The implementation therefore estimates $s$ from the two previous load steps
and caps the step at a fraction $f$ of the permissible value:
\begin{equation}\label{eq:eta-capped}
	\eta \leftarrow \min\left( \eta,\; \frac{f}{1-s} \right),
	\qquad
	s = \frac{\lambda_c - \lambda_c^{\,prev}}
	         {\lambda_b - \lambda_b^{\,prev}},
	\qquad f = 0.7
\end{equation}

Substituting Eq.~\ref{eq:eta-capped} into Eq.~\ref{eq:step-overshoot} gives
$\lambda_b' - \lambda_c' = -(1-f)\,(\lambda_c - \lambda_b)$ whatever the value
of $s$, so the cap does not merely prevent the overshoot: it turns
Eq.~\ref{eq:step} into a secant iteration on $\lambda_c(\lambda_b) -
\lambda_b$, which divides the remaining distance by a fixed factor at every
step.

The fraction is $0.7$ rather than something closer to unity because $s$ is
estimated backwards and lags while it steepens.  On the $n_y = 60$ mesh
the offending step is taken with $s = -0.66$ measured over the previous
interval, which permits $\eta = 0.54$ and therefore does not bind on the fixed
$\eta = 0.5$, while the step itself turns out to have $s = -1.09$.  A margin of
about a quarter on $1-s$ absorbs that lag, and with $f = 0.7$ the same mesh
lands at $\lambda_b/\lambda_c = 0.9965$ instead of $1.0013$, in the same number
of load steps.


\subsection{Accepting the state}

Convergence is declared by Sun~Eq.~(45).  That criterion measures the distance
to the bifurcation point from below only, so that once $\lambda_b$ passes
$\lambda_c$ its left-hand side turns negative and it accepts the state
whatever the overshoot.  The distance is therefore taken in absolute value:
\begin{equation}\label{eq:convergence}
	\frac{\left| \lambda_c-\lambda_b \right|}{\lambda_c} \le \varepsilon_1,
	\qquad \varepsilon_1 = 0.005
\end{equation}

If the load level lands past the bifurcation point by more than
$\varepsilon_1$, the load step is taken again from the previous state with
$\eta$ halved.  The fundamental path is still computable there, the
pre-buckling solve being restricted to the axisymmetric subspace, but the
expansion is meant to be made on the near side of it.

At convergence $\mu \simeq 1$ to within $\varepsilon_1$, which is what makes
the operators of Eq.~\ref{eq:eigenproblem} usable directly in the expansion.

\subsection{Rates of the pre-buckling state}

The first rate follows from Sun~Eq.~(29), using the tangent stiffness at the
converged pre-buckling state and solved in the reduced basis of
Eq.~\ref{eq:reduced-newton}:
\begin{equation}\label{eq:udot}
	\boldsymbol{K}_T\,\boldsymbol{\hat{u}}_{,\lambda} = \boldsymbol{F}_{ref}
\end{equation}

The second rate follows from Sun~Eq.~(35) as a backward difference against the
load step preceding the converged one:
\begin{equation}\label{eq:uddot}
	\boldsymbol{\hat{u}}_{,\lambda\lambda}
	= \frac{\boldsymbol{\hat{u}}_{,\lambda} - \boldsymbol{\hat{u}}_{,\lambda}^{\,prev}}
	       {\lambda_b - \lambda_b^{\,prev}}
\end{equation}

Reusing a step the load stepping has already converged on costs nothing and,
unlike a step solved ahead of $\lambda_b$, cannot fall on the far side of the
bifurcation point.  The corresponding stress rates, Sun~Eqs.~(30) and (34),
follow from Eqs.~\ref{eq:epsdot} and \ref{eq:epsddot} through the constitutive
matrix.

%==============================================================================
\section{Buckling modes of a cylinder}
\label{sec:canonical}
%==============================================================================

\subsection{The modes come in degenerate pairs}

For a circular cylinder the buckling modes occur in degenerate pairs, one for
each sign of the circumferential wave number, so $\dim\ker \ge 2$ even when a
single Koiter mode is retained, and every rotation of a pair inside its own
eigenspace is again a pair of buckling modes.  Which member of it the eigen
solver returns is decided by round-off, and changes with the BLAS
implementation.

That would not matter if the post-buckling coefficients were invariant under
that rotation, and they are not.  The second-order field of a mode with $n$
circumferential waves carries the harmonic $2n$, which a mesh sized for the
buckling mode itself barely resolves.  Measured on the case of
Section~\ref{sec:results} with $n_y = 40$ and $n = 10$, rotating the
critical pair by $30$ degrees moves $b_{1111}$ from $-0.230556$ to
$-0.192752$, sixteen per cent, with the pre-buckling state and the buckling
load unchanged to eleven digits.  Before the tangent of
Section~\ref{sec:tangent} was made consistent the same rotation moved it from
$-0.0008$ to $+0.2744$, sign included, so most of that dependence was the
tangent and what remains is the mesh.

\subsection{Fixing the member}

Each pair is rotated to the member whose crest falls on the $y = 0$
generator, which is a property of the mesh and not of the arithmetic.  Of the
two members of a pair one has a crest on that generator and the other a node,
so the dominant eigenvector of the $2\times2$ Gram matrix of their traces
there is well separated, and the sign is fixed by requiring the crest to point
outwards.

The partner of a mode need not be among the ones the eigen solver returned.  A
Krylov method builds its subspace from a single starting vector, so it finds
one direction of a degenerate eigenspace and then moves on to the next
multiplier: asking for the two lowest multipliers of the shell of
Section~\ref{sec:results} returns one member of the critical pair and one
member of the next pair, $3.3\times10^{-4}$ away.  A missing partner is
recovered through the cyclic shift $\boldsymbol{P}$ of
Section~\ref{sec:axisymmetric}, which maps a mode onto another mode of the
same pair, after removal of the axial rigid body translation that restores the
single-node constraint.  Truncating the offending degree-of-freedom instead
takes the mode out of its eigenspace and, the operator of the second-order
field being singular along it, the error is then amplified without bound.

This makes the coefficients reproducible: $9\times10^{-9}$ between a single
threaded and a multi threaded run, against a sign difference between operating
systems before.  It does not make them mesh converged, and only a mesh that
resolves the harmonic $2n$ would.

\subsection{How many modes the expansion actually needs}
\label{sec:howmany}

Fixing the member makes a single-mode result reproducible.  It does not make
it complete, and the measured spectra say by how much it falls short.  Twelve
multipliers were computed about the converged pre-buckling state of each
reference shell, $n_y = 40$; every one of them came out as a degenerate pair,
to eleven digits, so the cluster is counted in pairs below.

\begin{center}
\begin{tabular}{lcccc}
\toprule
 & \multicolumn{3}{c}{pairs within} & gap after \\
\cmidrule(lr){2-4}
 & $0.5\%$ & $1\%$ & $2\%$ & the $1\%$ cluster \\
\midrule
Sun~\S3.1, \texttt{NLprebuck=False}    & 1 & 3 & 5 & $0.28\%$ \\
Sun~\S3.1, \texttt{NLprebuck=True}     & 3 & 4 & 4 & $2.29\%$ \\
Arbocz, \texttt{NLprebuck=False}       & 1 & 4 & 6 & $0.26\%$ \\
Arbocz, \texttt{NLprebuck=True}        & 2 & 3 & 4 & $0.99\%$ \\
\bottomrule
\end{tabular}
\end{center}

Two readings of this table matter.

The first is that the single-mode expansion retains \emph{one} eigenvector out
of six or eight.  Its very first omission is the degenerate partner of the
critical mode, which Section~\ref{sec:canonical} goes to some trouble to
construct and then discards.  That alone bounds what the single-mode
coefficient can be expected to reproduce, and it is consistent with the
residual five per cent between $b_{1111}$ and the ANILISA value of
Section~\ref{sec:results} surviving every mesh refinement tried.  Closing the
$1\%$ cluster requires \texttt{koiter\_num\_modes} $= 6$ for the Arbocz shell
and $8$ for the Sun shell; closing only the $0.5\%$ cluster requires $4$ and
$6$.  In pairs, which is the natural unit here, that is three to four pairs.

The second is that the non-linear pre-buckling state does not merely move the
modes, it \emph{separates} them.  With \texttt{NLprebuck=False} the cluster
has no boundary worth the name, the next pair sitting $0.26$ to $0.28$ per
cent above the last one retained, so any truncation is arbitrary.  With
\texttt{NLprebuck=True} a genuine spectral gap opens, $0.99$ and $2.29$ per
cent, an order of magnitude wider than the spacing inside the cluster.  Only
then is ``the cluster'' a well-defined set, and only then is a multi-mode
truncation of it well posed.  The pre-buckling treatment of this document is
therefore a prerequisite for multi-mode Koiter, not an alternative to it.

Two pieces of the implementation stand between this and a working multi-mode
analysis.  The symmetrisation of $\phi_{30}$ over distinct mode indices is
established only for coinciding indices (Section~\ref{sec:limits}), and the
row border $\boldsymbol{W}$ of the bordered system
(Section~\ref{sec:secondorder}) carries one condition per retained Koiter
mode while the null space of $\phi_2$ has as many dimensions as the cluster
has members, so the two counts have to be brought into agreement before the
extra modes can be retained rather than deflated.

%==============================================================================
\section{Consistency of the operators}
\label{sec:consistency}
%==============================================================================

The operator appearing in the asymptotic expansion,
\begin{equation}\label{eq:phi2}
	\phi_2 = \boldsymbol{K}_C(\boldsymbol{\hat{u}})
	       + \mu\,\boldsymbol{K}_G(\boldsymbol{\hat{u}})
\end{equation}

\noindent must be \emph{the same} operator whose null vector is the buckling
mode, that is, exactly the operator of the eigenvalue problem of
Eq.~\ref{eq:eigenproblem}.  Building $\phi_2$ from a differently evaluated
$\boldsymbol{K}_C^{NL}$ than the one used there, for instance from
$\boldsymbol{K}_C^{NL}(\lambda_c\boldsymbol{\hat{u}})$ while the eigenvalue
problem used $\boldsymbol{K}_C^{NL}(\boldsymbol{\hat{u}})$, leaves $\phi_2$
non-singular along $\delta\boldsymbol{\bar{u}}$ and corrupts every quantity
derived from it.  The severity is easy to measure: with the inconsistent
pairing the expansion operator does not annihilate the buckling mode at all,
the relative residual rising from $10^{-10}$ to $4.6$.  Since
$\lambda_b\simeq\lambda_c$ after the iterative algorithm has converged,
$\mu\simeq1$ and Eq.~\ref{eq:phi2} is simply the tangent operator at the
critical state.

The discrete energy derivatives assembled element by element are the building
blocks of the expansion:

\begin{center}
\begin{tabular}{ll}
\toprule
symbol in the code & meaning \\
\midrule
\texttt{phi2}                & $\phi_2 = \partial^2\phi/\partial\boldsymbol{\bar{u}}^2$, the tangent operator at the critical state \\
\texttt{phi3\_ab[(i,j)]}     & $\phi_3$ contracted with modes $i,j$, one free index \\
\texttt{phi4[(i,j,k,l)]}     & $\phi_4$ fully contracted \\
\texttt{phi20\_a[i]}         & $\partial\phi_2/\partial\lambda$ contracted with mode $i$, one free index \\
\texttt{phi30\_ab[(i,j)]}    & $\partial\phi_3/\partial\lambda$ contracted with modes $i,j$, one free index \\
\texttt{phi200\_ab[(i,j)]}   & $\partial^2\phi_2/\partial\lambda^2$ fully contracted \\
\bottomrule
\end{tabular}
\end{center}

All of them are built from Eqs.~\ref{eq:epsdot} to \ref{eq:epsderivs} together
with the constant quadratic form
$\partial^2\boldsymbol{\varepsilon}/
\partial\boldsymbol{\bar{u}}\partial\boldsymbol{\bar{u}}$, which for
Eq.~\ref{eq:strain} is $\boldsymbol{G}_x^\top\boldsymbol{G}_x$,
$\boldsymbol{G}_y^\top\boldsymbol{G}_y$ and
$\boldsymbol{G}_x^\top\boldsymbol{G}_y +
\boldsymbol{G}_y^\top\boldsymbol{G}_x$ in the three strain components.

%==============================================================================
\section{Second-order fields and the orthogonality condition}
\label{sec:secondorder}
\label{sec:ortho}
%==============================================================================

\subsection{The singular system}

The second-order fields $\boldsymbol{\bar{u}}_{ij}$ solve
\begin{equation}\label{eq:secondorder}
	\phi_2\,\boldsymbol{\bar{u}}_{ij} = \boldsymbol{g}_{ij},
	\qquad
	\boldsymbol{g}_{ij} = -\frac{1}{2}\,\phi_{3,ij}
	  - \frac{1}{m}\sum_{k} a_{kij}\,\lambda_k\,\phi_{20,k}
\end{equation}

\noindent with $m$ the number of Koiter modes.  By construction $\phi_2$ is
singular, its null space being spanned by the buckling modes.  Consequently
Eq.~\ref{eq:secondorder} determines $\boldsymbol{\bar{u}}_{ij}$ only up to an
arbitrary element of that null space, and a plain sparse solve followed by a
projection is not a well-posed procedure.  Two things are therefore needed: a
solver that copes with the singularity, and the correct condition that fixes
the null-space component.

\subsection{The condition to be imposed}

The orthogonality condition of Koiter's theory is Sun~Eq.~(18), whose finite
element form is Sun~Eq.~(33):
\begin{equation}\label{eq:sun33}
	\delta\boldsymbol{\bar{u}}_k^\top
	\left[ \boldsymbol{K}_D(\boldsymbol{\hat{u}},\boldsymbol{\hat{u}}_{,\lambda})
	+ \boldsymbol{K}_G(\boldsymbol{\sigma}_{,\lambda}) \right]
	\boldsymbol{\bar{u}}_{ij}
	+ \frac{1}{2}\,
	\delta\boldsymbol{\bar{u}}_k^\top
	\boldsymbol{B}_{NL}^\top(\delta\boldsymbol{\bar{u}}_k)\,\boldsymbol{C}\,
	\boldsymbol{B}_{NL}(\delta\boldsymbol{\bar{u}}_k)\,
	\boldsymbol{\hat{u}}_{,\lambda} = 0
\end{equation}

Two observations turn this into something directly computable.  The bracket in
the first term is the load-parameter derivative of the tangent operator,
$\partial \phi_2/\partial\lambda$, so the first term is exactly the
contraction stored as \texttt{phi20\_a[k]}.  The second term is a constant,
which follows from $\phi_{30}$ without any new element integral: $\phi_{30,ij}$
is assembled as one half of the quadrature weight times the sum of six
contractions which, when all mode indices coincide, are equal to one another
by the symmetry of the constitutive matrix and of the quadratic form.  Hence
\begin{equation}\label{eq:cconst}
	c_k = \frac{\phi_{30,ij}\cdot\delta\boldsymbol{\bar{u}}_k}{6}
\end{equation}

\noindent and the condition to impose is
\begin{equation}\label{eq:condition}
	\phi_{20,k}\cdot\boldsymbol{\bar{u}}_{ij} = -\,c_k
\end{equation}

It is weighted by the stiffness matrices, through $\phi_{20}$, and
inhomogeneous, through $c_k$.  A Gram-Schmidt projection, or equivalently a
symmetric Euclidean border, imposes neither of these properties and is the
procedure that Eq.~\ref{eq:condition} replaces.

\subsection{The bordered system}

Both requirements are met by a \emph{non-symmetric} bordered system:
\begin{equation}\label{eq:bordered}
	\begin{bmatrix}
	  \phi_2 & \boldsymbol{N}_{sp} \\[2pt]
	  \boldsymbol{W}^\top & \boldsymbol{0}
	\end{bmatrix}
	\begin{Bmatrix} \boldsymbol{\bar{u}}_{ij} \\ \boldsymbol{\alpha} \end{Bmatrix}
	=
	\begin{Bmatrix} \boldsymbol{g}_{ij} \\ -\boldsymbol{c} \end{Bmatrix}
\end{equation}

\noindent in which the roles of the two borders are distinct.  The
\emph{column} border $\boldsymbol{N}_{sp}$ spans the null space of $\phi_2$ and
is what makes the system non-singular; the \emph{row} border $\boldsymbol{W}$
carries the constraints, one column per Koiter mode, $\boldsymbol{W}[:,k] =
\phi_{20,k}$, with the corresponding right-hand side $-c_k$ from
Eq.~\ref{eq:cconst}.  The system is non-singular provided
$\boldsymbol{W}^\top\boldsymbol{N}_{sp}$ is, and its leading entry
$\phi_{20,k}\cdot\delta\boldsymbol{\bar{u}}_k$ is exactly the quantity
appearing in the denominator of the post-buckling coefficients, hence non-zero
by construction.

Because the critical eigenspace is degenerate, $\boldsymbol{N}_{sp}$ must span
the \emph{whole} near-null space and not only the retained Koiter modes: every
eigenvector whose multiplier equals the critical one within a relative
tolerance is included, and the set is orthonormalised with a rank-revealing
factorisation.  Equation~\ref{eq:sun33} supplies one condition per Koiter
mode, which is fewer than the dimension of the null space; the degenerate
partners lie outside single-mode theory, and the remaining rows of
$\boldsymbol{W}$ keep the Euclidean condition.  This is a choice, not a
result; the consistent treatment of a degenerate critical mode is the
multi-mode Koiter expansion.

\subsection{What the weighted condition does, and does not, change}

Condition~\ref{eq:condition} is the correct one and replaces a Gram-Schmidt
projection that was neither weighted nor inhomogeneous.  For the cases
considered here, however, it changes almost nothing.  Measured before the
tangent was made consistent, running the same analysis under the two
conditions with everything else identical moved the second-order field by
$6.5\times10^{-7}$ relative, so the two conditions are numerically
near-coincident for this problem: the null-space component of
$\boldsymbol{\bar{u}}_{ij}$ on which they disagree is small to begin with.

A second reason reinforces this.  Condition~\ref{eq:condition} fixes only the
component of $\boldsymbol{\bar{u}}_{ij}$ along $\ker\phi_2$, and $b_{ijkl}$
sees $\boldsymbol{\bar{u}}_{ij}$ through the contractions
$\phi_{3,ij}\cdot\boldsymbol{\bar{u}}_{kl}$; but $\phi_{3,ij}$ contracted with
a buckling mode is precisely the numerator of $a_{ijk}$ in
Eq.~\ref{eq:a-coefficient}.  For the symmetric bifurcation of a cylinder under
axial compression $a_{ijk}$ vanishes, measured values being of order
$10^{-5}$ against a $b_{ijkl}$ of order unity, so $\phi_{3,ij}$ is nearly
orthogonal to the null space.  For an asymmetric bifurcation, with
$a_{ijk}\neq0$, the condition would carry more weight.

%==============================================================================
\section{Post-buckling coefficients}
%==============================================================================

With the second-order fields available, the coefficients are
\begin{equation}\label{eq:a-coefficient}
	a_{ijk} = -\frac{1}{2\lambda_i}\,
	  \frac{\phi_{3,ij}\cdot\delta\boldsymbol{\bar{u}}_k}
	       {\phi_{20,i}\cdot\delta\boldsymbol{\bar{u}}_i}
\end{equation}
\begin{multline}\label{eq:b-coefficient}
	b_{ijkl} = -\frac{1}
	  {6\,\lambda_i\,(\phi_{20,i}\cdot\delta\boldsymbol{\bar{u}}_i)}
	\Big[\,
	  \phi_{4,ijkl}
	  + 3\,\phi_{3,ij}\cdot\boldsymbol{\bar{u}}_{kl}
	  + 3\,\phi_{3,il}\cdot\boldsymbol{\bar{u}}_{jk} \\
	  + \lambda_i\big(
	      a_{iij}\,\phi_{30,ik}\!\cdot\!\delta\boldsymbol{\bar{u}}_l
	    + a_{ijk}\,\phi_{30,il}\!\cdot\!\delta\boldsymbol{\bar{u}}_i
	    + a_{ikl}\,\phi_{30,ii}\!\cdot\!\delta\boldsymbol{\bar{u}}_j \big) \\
	  + \phi_{200,ii}\,\lambda_i^2\big(
	      a_{iij}a_{ikl} + a_{ijk}a_{ili} + a_{ikl}a_{iij} \big)
	\Big]
\end{multline}

The buckling modes are normalised so that the largest translational amplitude
equals the shell thickness, which is the normalisation used by the reference
solutions.

%==============================================================================
\section{Algorithm}
%==============================================================================

\begin{algorithm}[H]
\caption{Koiter analysis about a non-linear pre-buckling state}
\begin{algorithmic}[1]
\State solve $\boldsymbol{K}_0\boldsymbol{\bar{u}} = \boldsymbol{F}_{ref}$;
       project onto the axisymmetric subspace; $\lambda_b \gets 1$
\For{$i = 1,\dots,i_{\max}$}
  \State assemble $\boldsymbol{K}_C(\boldsymbol{\hat{u}})$,
         $\boldsymbol{K}_G(\boldsymbol{\hat{u}})$
  \State solve Eq.~\ref{eq:eigenproblem} $\rightarrow \mu$,
         $\delta\boldsymbol{\bar{u}}$; canonicalise the degenerate pairs
         (Section~\ref{sec:canonical}); \quad $\lambda_c \gets \lambda_b\mu$
  \State \textbf{if} Eq.~\ref{eq:convergence} holds \textbf{then break}
  \State \textbf{if} $\lambda_b > \lambda_c$ \textbf{then} restore the
         previous step, halve the cap on $\eta$, \textbf{continue}
  \State $\eta \gets 0.8$ if $i=1$ else $0.5$, capped by
         Eq.~\ref{eq:eta-capped}
  \Repeat
    \State $\lambda_b^{new} \gets \lambda_b + \eta(\lambda_c-\lambda_b)$
    \State predictor $\boldsymbol{\hat{u}} +
           \boldsymbol{\hat{u}}_{,\lambda}(\lambda_b^{new}-\lambda_b)$;
           Newton-Raphson in the reduced basis
           (Eq.~\ref{eq:reduced-newton})
    \State \textbf{on failure} $\eta \gets \eta/2$
  \Until{converged or $\eta$ exhausted}
  \State store $(\lambda_b,\lambda_c,\boldsymbol{\hat{u}}_{,\lambda})$ as the
         previous step; advance
\EndFor
\State $\boldsymbol{\hat{u}}_{,\lambda}$ from Eq.~\ref{eq:udot};
       $\boldsymbol{\hat{u}}_{,\lambda\lambda}$ from Eq.~\ref{eq:uddot}
\State assemble $\phi_2,\phi_3,\phi_4,\phi_{20},\phi_{30},\phi_{200}$
\State $a_{ijk}$ from Eq.~\ref{eq:a-coefficient}
\State $\boldsymbol{\bar{u}}_{ij}$ from the bordered system
       Eq.~\ref{eq:bordered} with the weighted condition
       Eq.~\ref{eq:condition}
\State $b_{ijkl}$ from Eq.~\ref{eq:b-coefficient}
\end{algorithmic}
\end{algorithm}

%==============================================================================
\section{Effect on the results}
\label{sec:results}
%==============================================================================

The two verification cases are the composite cylinder of Sun~\S3.1 and the
NASA shell AW-CYL-1-1 of \citet{arbocz2001}; they are
the same shell, so the ANILISA $n$-search of the latter provides reference
values for both.  Buckling loads are normalised by $N_{c\ell} = E_{11}h^2 /
\left(r\sqrt{3(1-\nu_{12}^2)}\right)$.

\subsection{Buckling loads}

\begin{center}
\begin{tabular}{lccc}
\toprule
 & $n$ & $\lambda_c$ & reference \\
\midrule
Arbocz, $n_y = 40$ & $10$ & $0.331413$ & $0.329163$ (ANILISA, $n=10$) \\
Arbocz, $n_y = 60$ & $11$ & $0.330603$ & $0.328594$ (ANILISA, $n=11$) \\
                        &      &            & $0.327759$ (STAGS-A, $n=11$) \\
Sun, $n_y = 40$    & $10$ & $0.346127$ & --- \\
Sun, $n_y = 50$    & $10$ & $0.347244$ & --- \\
\bottomrule
\end{tabular}
\end{center}

The two Arbocz meshes agree with the ANILISA entry for their own
circumferential wave number to $0.7$ and $0.6$ per cent respectively, and
the finer one with the Level-3 STAGS-A value to $0.9$ per cent.  The membrane
pre-buckling regime, which is what the computation returns when the expansion
point is left at the reference load, gives $0.365992$ instead.

\subsection{Post-buckling coefficients}

\begin{center}
\begin{tabular}{lccc}
\toprule
 & $n$ & $b_{1111}$ & reference \\
\midrule
Arbocz, $n_y = 40$ & $10$ & $-0.335975$ & --- \\
Arbocz, $n_y = 60$ & $11$ & $-0.356457$ & $-0.37605$ (ANILISA, $n=11$) \\
                        &      &             & $-0.3772$ (Sun, Table~2) \\
Sun, $n_y = 40$    & $10$ & $-0.230556$ & --- \\
Sun, $n_y = 50$    & $10$ & $-0.222844$ & --- \\
\bottomrule
\end{tabular}
\end{center}

The $n_y = 60$ value is within five per cent of the ANILISA coefficient
for the same mode.  A negative $b$ means an imperfection sensitive shell,
which is what this one is.

\subsection{What the tangent correction changed}

Before the tangent stiffness matrix was made consistent
(Section~\ref{sec:tangent-check}) the same meshes returned $\lambda_c =
0.302887$ and $b_{1111} = +0.410664$ for the Arbocz case: the buckling load
was nine per cent low and the post-buckling coefficient had the wrong sign.
The sign was at that time attributed to a single-mode expansion about a
degenerate critical mode being unable to determine it, which was incorrect.
Two meshes that did not even select the same critical mode, $n = 10$ at
$n_y = 40$ against $n = 7$ at $n_y = 50$ for the Sun case, now
select $n = 10$ on both and agree to $0.3$ per cent in the buckling load and
$3.4$ per cent in $b_{1111}$.

The Newton-Raphson behaviour changed with it.  Each load step now takes a
single iteration, with residuals falling $5.1\times10^{-5}$,
$1.1\times10^{-5}$, $1.2\times10^{-6}$, $5.4\times10^{-8}$,
$1.5\times10^{-9}$, and no step back-off anywhere.  The $n_y = 60$ mesh,
which previously exhausted the iteration budget at every load step, converges
in four of them.

%==============================================================================
\section{The variable-stiffness CTS model}
\label{sec:cts}
%==============================================================================

The routines \texttt{fkoiter\_cylinder\_CTS\_circum} of
\texttt{koiter\_cylinder\_CTS.py} and \texttt{koiter\_cylinder\_CTS\_sanders.py}
apply the whole of the preceding development to the continuous-tow-shearing
(CTS) variable-stiffness cylinders of dos~Santos and Castro,
\emph{Lightweight design of variable-stiffness cylinders with reduced
imperfection sensitivity enabled by continuous tow shearing and machine
learning}, Materials \textbf{15} (2022) 4117.  The algorithm is taken
literally from the constant-stiffness models, and a check in the test suite
asserts that the shared regions of the two files are identical character for
character, so only what is specific to the CTS model is documented here.

\subsection{Why the axisymmetric reduction still applies}
\label{sec:cts-axi}

The parameterisation steers the tows along the \emph{axial} direction only,
$\theta = \theta(x)$, and the tow shearing thickens the laminate as
$h_{tow}/\cos\theta$, again a function of $x$ alone.  The stiffness is
therefore uniform around the circumference, the cyclic shift of
Section~\ref{sec:axisymmetric} remains a symmetry of every assembled operator,
and the fundamental path of the perfect cylinder under uniform axial
compression is still axisymmetric.  The Galerkin reduction of
Eq.~\ref{eq:reduced-newton} carries over unchanged.

Two consequences are worth stating explicitly, because they are easy to
confuse:

\begin{enumerate}\itemsep3pt
  \item The reduction constrains the \textbf{pre-buckling state only}.  The
        eigenvalue problem is assembled and solved on the full space, so
        non-axisymmetric buckling modes are obtained exactly as they are in
        the constant-stiffness models.  What the reduction removes is the
        possibility of the pre-buckling iteration drifting off the
        fundamental path, not the modes.
  \item The \emph{premise} of the reduction --- that every buckling mode
        making $\boldsymbol{K}_T$ singular is non-axisymmetric --- is
        \textbf{not} satisfied everywhere in the CTS design space.  See
        Section~\ref{sec:cts-spectrum}.
\end{enumerate}

Unlike the constant-stiffness models, the axial station spacing is not
uniform: it is built from the transition length $t = r_{CTS}\sin|\theta_{c2} -
\theta_{c1}|$ and the plateau lengths $c_1$ and $c_2$.  An assertion checks
that the stations \texttt{mesh\_order} sorts the nodes into are the ones the
axisymmetric basis is then built on, which a non-uniform $x$ grid could
otherwise violate silently.

\subsection{Length scales that size the mesh}
\label{sec:cts-scales}

Three scales control the discretisation.  With $r = 68$~mm, $L = 300$~mm and
$h \simeq 0.8$~mm, the geometry of the cylinder used throughout the test
suite, they are

\begin{center}
\begin{tabular}{lcl}
\toprule
scale & value & constrains \\
\midrule
$\sqrt{rh}$ & $7.4$~mm & --- \\
edge boundary layer $1/\beta = \sqrt{rh}\,\big/\,[3(1-\nu^2)]^{1/4}$
    & $5.8$~mm & $n_x$ \\
buckling half-wave $\pi\sqrt{rh}\,\big/\,[12(1-\nu^2)]^{1/4}$
    & $12.8$~mm & $n_x$ \\
wavelength of the $2n$ harmonic, $\pi r/n$ & $21$~mm at $n{=}10$ & $n_y$ \\
\bottomrule
\end{tabular}
\end{center}

The axial mesh is set by the edge boundary layer.  With
\texttt{NLprebuck=True} the non-linear pre-buckling deformation caused by the
edge restraint is concentrated within a few multiples of $1/\beta$ of each
end, and it is precisely that deformation which the expansion is about; with
\texttt{NLprebuck=False} the requirement is much weaker, the linear
pre-buckling state being smooth.  This is a demand the earlier
implementation did not make, because expanding at the reference load left the
boundary layer undeveloped.

The circumferential mesh is set by the \emph{second-order field}, not by the
buckling mode.  A mode with $n$ circumferential waves drives a second-order
field carrying the harmonic $2n$ (Section~\ref{sec:limits}), so a mesh sized
to resolve the mode itself samples the field that determines $b_{ijkl}$ at
only about two nodes per wave.  Resolving the $2n$ harmonic with four nodes
per wave requires $n_y \simeq 8n$.

\subsection{Measured convergence}
\label{sec:cts-convergence}

The design below is representative of the parameterisation: $\theta_{c1} =
30^\circ$, $\theta_{c2} = 60^\circ$, $n = 3$, $c_{2,ratio} = 0.5$, $r_{CTS} =
50$~mm, giving $t = 25$~mm, $c_1 = 18.75$~mm and $c_2 = 25$~mm.  Sanders
kinematics, \texttt{idealistic\_CTS=True}, \texttt{NLprebuck=False}, four
integration points.  The axial mesh is refined with the circumferential one
through \texttt{nxt}, which keeps the element aspect ratio between $1.0$ and
$1.4$.

\begin{center}
\begin{tabular}{rrrrrrr}
\toprule
$n_y$ & $n_x$ & DOF & $\Delta x$ & $\Delta x\,\beta$ & $P_{cr}$~[N] &
$\Delta P_{cr}$ \\
\midrule
 30 & 23 &  6\,900 & $13.64$~mm & $2.37$ & $40982.6$ & $+6.65\%$ \\
 45 & 37 & 16\,650 & $ 8.33$~mm & $1.45$ & $38893.0$ & $+1.21\%$ \\
 60 & 49 & 29\,400 & $ 6.25$~mm & $1.09$ & $38572.9$ & $+0.38\%$ \\
 80 & 71 & 56\,800 & $ 4.29$~mm & $0.74$ & $38452.1$ & $+0.06\%$ \\
100 & 97 & 97\,000 & $ 3.12$~mm & $0.54$ & $38427.6$ & --- \\
\bottomrule
\end{tabular}
\end{center}

The buckling load converges monotonically, and it does so as a function of
$\Delta x\beta$: the error falls below one per cent once the axial element is
about as long as the boundary layer, $\Delta x \simeq 1/\beta$, and below one
tenth of a per cent at $\Delta x \simeq 0.75/\beta$.  The run times on one
core were $24$, $65$, $91$, $150$ and $323$~seconds.

The post-buckling coefficient behaves completely differently.  On the same
five meshes $b_{1111}$ was
\[
  -2.6\times10^{-3},\quad
  +1.2\times10^{-2},\quad
  +1.9\times10^{-3},\quad
  +3.7\times10^{-4},\quad
  +5.0\times10^{-5}
\]
changing sign and two orders of magnitude.  This is not a mesh resolution
problem and no refinement repairs it; the cause is the buckling spectrum of
this particular design, treated next.

\subsection{The same design with a non-linear pre-buckling state}
\label{sec:cts-nl}

Repeating the study with \texttt{NLprebuck=True} changes both numbers and
behaviour:

\begin{center}
\begin{tabular}{rrrrrrr}
\toprule
$n_y$ & $n_x$ & $P_{cr}$~[N] & $\Delta P_{cr}$ & $n$ & $b_{1111}$ &
$\Delta b_{1111}$ \\
\midrule
30 & 23 & $35087.8$ & $+24.59\%$ & $4$ & $-0.0565$ & $-62.7\%$ \\
45 & 37 & $28517.4$ & $+1.26\%$  & $7$ & $-0.1791$ & $+18.4\%$ \\
60 & 49 & $28309.7$ & $+0.52\%$  & $7$ & $-0.1564$ & $+3.3\%$ \\
80 & 71 & $28161.9$ & ---        & $7$ & $-0.1514$ & --- \\
\bottomrule
\end{tabular}
\end{center}

Three things change for the better.  The buckling load is $27$ per cent below
the linear-pre-buckling value of Section~\ref{sec:cts-convergence}, as the
edge restraint would lead one to expect.  The critical wave number settles at
$n = 7$ from $n_y = 45$ upwards instead of wandering.  And $b_{1111}$
\emph{converges}: $-0.179$, $-0.156$, $-0.151$, monotonically and without
changing sign, to within $3.3$ per cent at $n_y = 60$, where the linear
analysis of Section~\ref{sec:cts-convergence} produced two orders of
magnitude of sign-changing noise on the same meshes.  The post-buckling
coefficient of a CTS cylinder is therefore a computable quantity, but only
about a non-linear pre-buckling state.

The coarsest mesh is useless here in a way it was not for the linear
analysis: $n_y = 30$ is $25$ per cent off in $P_{cr}$ and selects a different
mode.  That is the boundary-layer requirement of
Section~\ref{sec:cts-scales} asserting itself, the expansion point now living
where the edge restraint acts.

\medskip
\noindent
\textbf{The default iteration budget is not enough.}  With the default
\texttt{NLprebuck\_maxiter=12} the $n_y = 45$ run exhausted the budget and
stopped at $\lambda_b/\lambda_c = 0.9933$, short of the $0.995$ requested,
after one Newton-Raphson back-off and two overshoot-and-step-back cycles;
the table above was produced with \texttt{NLprebuck\_maxiter=30}, which
converged every mesh with no warning at all.  The default was tuned on the
constant-stiffness shells of Section~\ref{sec:results}, where $\lambda_c$
moves less with $\lambda_b$ and four to seven load steps suffice.  A CTS
design with a dense spectrum spends steps backing off, and the budget has to
be raised accordingly.  The symptom is explicit --- the run prints
\texttt{WARNING: the iterative eigenvalue algorithm stopped at
lambda\_b/lambda\_c = ...} --- and must not be ignored, because a state
stopped short of the requested neighbourhood is exactly the regime
Section~\ref{sec:step} exists to avoid.

\subsection{The buckling spectrum is design-dependent}
\label{sec:cts-spectrum}

The circumferential harmonic content of the lowest modes differs
qualitatively across the design space.  Measured on the $n_y = 60$ mesh with
a \emph{linear} pre-buckling state:

\begin{center}
\begin{tabular}{lccl}
\toprule
design & critical mode & $\mu_1/\mu_0$ & gap to the next distinct mode \\
\midrule
$\theta_{c1}=\theta_{c2}=45^\circ$ & axisymmetric, $n=0$ & $1.00018$ & $0.63\%$ \\
$\theta_{c1}=30^\circ$, $\theta_{c2}=60^\circ$, $n{=}3$
                                   & axisymmetric, $n=0$ & $1.00344$ & $0.34\%$ \\
$\theta_{c1}=\theta_{c2}=60^\circ$ & pair, $n=3$         & $1.00000$ & $14.19\%$ \\
\bottomrule
\end{tabular}
\end{center}

The $\pm60^\circ$ design is the textbook situation the single-mode expansion
assumes: a clean degenerate pair, well separated from everything above it, and
with $n = 3$ the $2n = 6$ harmonic is sampled by ten nodes per wave already at
$n_y = 60$.

The other two are not.  Their critical mode is \emph{axisymmetric}, and the
first six multipliers lie within one per cent of each other, a dense cluster
rather than an isolated pair.  The single-mode $b_{ijkl}$ is then not
determinate, the expansion being made about one member of a cluster the mesh
reorders from one refinement to the next, and that is exactly the
oscillation of Section~\ref{sec:cts-convergence}.  Worse, the premise of the
axisymmetric reduction is violated: the critical mode lies \emph{inside} the
subspace the pre-buckling problem is solved in.

\medskip
\noindent
\textbf{The non-linear pre-buckling state removes the pathology.}  Repeating
the steered design with \texttt{NLprebuck=True} moves the critical mode off
the axisymmetric harmonic entirely: it comes out at $n = 4$ on the
$n_y = 30$ mesh and $n = 7$ on both $n_y = 45$ and $n_y = 60$, and
$b_{1111}$ becomes $-0.179$ and $-0.156$ on the latter two, values of the
same order as the reference shells of Section~\ref{sec:results} rather than
the near-zero noise of the linear case.  The axisymmetric critical mode is
therefore an artefact of expanding about the \emph{linear} pre-buckling
state, in which the edge restraint has not yet developed, and not a property
of the design.  This is the same separation of the spectrum documented in
Section~\ref{sec:howmany} for the constant-stiffness shells, seen here in its
sharpest form.

The practical consequence is a warning about screening.  A spectrum computed
with \texttt{NLprebuck=False} does not predict the one computed with
\texttt{NLprebuck=True}, neither the wave number nor the clustering, so the
check recommended in Section~\ref{sec:cts-recommendation} has to be run with
the same setting as the production analysis.  And if a design genuinely does
retain an axisymmetric critical mode under \texttt{NLprebuck=True}, the
failure is silent: the load stepping halts just short of the bifurcation
point, where the reduced tangent is still regular, and reports convergence
with no indication that the state it stopped at is about to bifurcate
\emph{within} the subspace it was computed in.  The diagnostic is the
harmonic content of the critical mode, never the convergence history.

\subsection{Recommendation}
\label{sec:cts-recommendation}

\begin{enumerate}\itemsep4pt
  \item \textbf{Always inspect the spectrum before using $b_{ijkl}$.}
        Run once with \texttt{koiter\_num\_modes=0}, which returns
        immediately after the eigenvalue analysis and costs a small fraction
        of a full run, and check (i) that the critical mode is a
        non-axisymmetric degenerate pair and (ii) that the next distinct
        multiplier is separated by appreciably more than
        \texttt{NLprebuck\_eps1}.  If either fails, the buckling load is
        still meaningful but the single-mode $b_{ijkl}$ is not, whatever the
        mesh.  Run this screening with the \emph{same} \texttt{NLprebuck}
        setting as the production analysis: Section~\ref{sec:cts-spectrum}
        shows the two give different wave numbers and different clustering on
        the same design.
  \item \textbf{For the buckling load}, size the axial mesh by the boundary
        layer: $\Delta x \le 1/\beta$ for about half a per cent, $\Delta x
        \simeq 0.75/\beta$ for about a tenth.  For the test-suite geometry
        that is $n_y = 60$ and $n_y = 80$ respectively, with
        \texttt{nxt} chosen to hold the aspect ratio near one.  Use the
        finer of the two with \texttt{NLprebuck=True}, where the boundary
        layer carries the expansion point.
  \item \textbf{For $b_{ijkl}$, when the spectrum permits it at all}, take
        $n_y \simeq 8n$ with $n$ the measured circumferential wave
        number, so that the $2n$ harmonic gets four nodes per wave.  For the
        $\pm60^\circ$ design, $n = 3$, this is satisfied from $n_y = 24$
        upwards, and the mesh is then governed by the buckling load
        requirement instead.
  \item \textbf{Do not transfer the $240$ elements of the reference.}  The
        $240$ S4R elements around the circumference used for the Abaqus
        cross-check in the reference are a linear reduced-integration shell
        requirement.  The BFSC element is $C^1$ with cubic interpolation and
        ten degrees of freedom per node and resolves the same harmonic with
        far fewer elements; $n_y = 60$ to $80$ is the corresponding
        range here.
  \item \textbf{Raise \texttt{NLprebuck\_maxiter}.}  The default $12$ is
        a constant-stiffness figure.  Use $30$ for CTS designs, and treat the
        \texttt{WARNING} about stopping short of
        $1 - \texttt{NLprebuck\_eps1}$ as a failure rather than a note: the
        expansion is then made about a state that is not in the neighbourhood
        of the bifurcation point, which is the very thing
        Section~\ref{sec:step} is for.
  \item \textbf{Cost.}  A full run at $n_y = 80$ takes of the order of
        minutes with \texttt{NLprebuck=False} and several times that with
        \texttt{NLprebuck=True}, the load stepping solving an eigenvalue
        problem per step.  In a design-of-experiments loop over thousands of
        samples, as the reference performs, this is the binding constraint,
        and the screening run of item~1 is the part worth doing on every
        sample.
\end{enumerate}

%==============================================================================
\section{Known limitations}
\label{sec:limits}
%==============================================================================

\begin{enumerate}\itemsep3pt
  \item \textbf{Degenerate critical modes.}  Condition~\ref{eq:condition}
        supplies one equation per retained mode while the null space is
        larger, and the rotation inside the degenerate pair still moves
        $b_{1111}$ by sixteen per cent on the $n_y = 40$ mesh.  The
        canonicalisation of Section~\ref{sec:canonical} makes the choice
        reproducible rather than arbitrary; it does not make it converged.
        Multi-mode Koiter is the consistent treatment, and
        Section~\ref{sec:howmany} sizes it: six to eight modes, three to four
        degenerate pairs, for the reference shells.
  \item \textbf{Resolution of the second-order field.}  The second-order field
        carries the harmonic $2n$, which the meshes used here sample with
        about two nodes per wave.  This, and not the expansion itself, is what
        limits $b_{ijkl}$; the buckling load is far less affected.
  \item \textbf{Symmetrisation of $\phi_{30}$ for distinct modes.}  The
        identity of Eq.~\ref{eq:cconst} was established for coinciding mode
        indices, which covers the single-mode case exercised here; the general
        multi-mode symmetrisation remains to be worked out.
  \item \textbf{Axisymmetric critical modes.}  The reduction of the
        pre-buckling problem to the axisymmetric subspace
        (Section~\ref{sec:axisymmetric}) assumes that no buckling mode lies
        in that subspace.  With a \emph{linear} pre-buckling state parts of
        the CTS design space violate it, the critical mode there being
        axisymmetric (Section~\ref{sec:cts-spectrum}).  On every design
        tried, switching to \texttt{NLprebuck=True} moved the critical mode
        off the axisymmetric harmonic and restored the assumption, so the
        violation has not been observed where the reduction is actually
        relied upon --- but it has not been \emph{excluded} either, and it
        would fail silently, the load stepping halting just short of the
        bifurcation point where the reduced tangent is still regular and
        reporting convergence.  No part of the code inspects the harmonic
        content of the critical mode, which is the only reliable
        diagnostic.
  \item \textbf{Load control.}  The scheme of Section~\ref{sec:step} now
        reaches the requested neighbourhood of the bifurcation point on all
        meshes tried, but it remains load controlled.  The extrapolated
        eigenvalue problem of Sun~Eq.~(27),
        $\left[\boldsymbol{K}_T + (\lambda_c-\lambda_b)\,
        \partial\boldsymbol{K}_T/\partial\lambda\right]
        \delta\boldsymbol{\bar{u}}=\boldsymbol{0}$, would remove the need to
        sit on top of the bifurcation point at all.
  \item \textbf{Quadrature.}  Convergence with respect to the number of Gauss
        points has been verified: $n_{int}=4$ gives $b$ to within $0.4$ per
        cent of $n_{int}=5$ and $6$, while $n_{int}=2$ under-integrates and
        produces a rank-deficient element.
\end{enumerate}

%==============================================================================
\appendix
\section{Reproducing the numbers}
\label{sec:appendix}
%==============================================================================

Every figure in this document that is not a literature value is produced by a
script kept in the repository, so that the document and the code cannot drift
apart.  Two kinds of script are involved.  The \emph{verification studies} of
Appendix~\ref{sec:appendix-scripts} sweep meshes and report; they take minutes
to tens of minutes and assert nothing.  The \emph{tests} of
Appendix~\ref{sec:appendix-tests} assert, run under \texttt{pytest}, and are
the ones that fail if a change breaks something.

All of them require \texttt{bfsccylinder} $\geq 0.6.0$, whose tangent
stiffness matrix is consistent with the internal force vector; several of the
figures quoted here measure exactly that consistency, and every one of them
changes on earlier versions (Section~\ref{sec:results}).

\subsection{Verification studies}
\label{sec:appendix-scripts}

In \texttt{doc/verification/}.  Each puts the repository root first on
\texttt{sys.path}, so an installed copy of the package cannot shadow the
working tree, and each is run directly, for instance
\begin{center}
\texttt{python doc/verification/kinematics\_vs\_element.py}
\end{center}

\begin{center}
\begin{tabular}{lll}
\toprule
script & produces & runtime \\
\midrule
\texttt{kinematics\_vs\_element.py}
  & Section~\ref{sec:consistency} & seconds \\
\texttt{tangent\_consistency.py}
  & Section~\ref{sec:tangent} & seconds \\
\texttt{cts\_vs\_constant\_stiffness.py}
  & Section~\ref{sec:cts} & $\sim$10~min \\
\texttt{cts\_mesh\_convergence.py}
  & table of Section~\ref{sec:cts-convergence} & $\sim$12~min \\
\texttt{cts\_mesh\_convergence.py nl}
  & table of Section~\ref{sec:cts-nl} & $\sim$35~min \\
\texttt{cts\_mode\_spectrum.py}
  & table of Section~\ref{sec:cts-spectrum} & $\sim$20~min \\
\texttt{reference\_mode\_cluster.py}
  & table of Section~\ref{sec:howmany} & $\sim$5~min \\
\bottomrule
\end{tabular}
\end{center}

\noindent
Two of these deserve a note.
\texttt{kinematics\_vs\_element.py} rebuilds $\boldsymbol{f}_{int}$ and the
second variation from the strain measures the Koiter expansion assumes and
compares them against the compiled element; it includes a deliberately wrong
control, the Sanders element driven with von Karman kinematics, which fails at
$7.8\times10^{-4}$ and so establishes that the agreement at $10^{-16}$ is not
vacuous.  \texttt{cts\_mesh\_convergence.py} runs with
\texttt{NLprebuck\_maxiter=30} rather than the default, for the reason given
in Section~\ref{sec:cts-nl}.

\subsection{Literature verification cases}
\label{sec:appendix-tests}

In \texttt{tests/}, run with \texttt{pytest} from that directory.  The
reference values and the meshes they belong to are documented in the tests
themselves; the regression values are for their mesh and are not converged,
which the comments there state case by case.

\begin{center}
\begin{tabular}{lll}
\toprule
test & shell & reference \\
\midrule
\texttt{test\_koiter\_cylinder\_newton\_raphson.py}
  & Sun~\S3.1; AW-CYL-1-1
  & \citet{sun2020}; \citet{arbocz2001} \\
\texttt{test\_koiter\_cylinder\_Waters.py}
  & Waters shell
  & Arbocz and Starnes (2002) \\
\texttt{test\_koiter\_cylinder\_Waters\_sanders.py}
  & Waters shell
  & Arbocz and Starnes (2002) \\
\texttt{test\_koiter\_cylinder\_CTS.py}
  & CTS, constant-stiffness limit
  & \citet{dossantos2022} \\
\texttt{test\_koiter\_cylinder\_CTS\_sanders.py}
  & CTS, constant-stiffness limit
  & \citet{dossantos2022} \\
\texttt{test\_buckling\_mode\_cluster.py}
  & Sun~\S3.1
  & Section~\ref{sec:howmany} \\
\texttt{test\_linBuck\_VAFW.py}
  & VAFW cylinders
  & linear buckling \\
\texttt{test\_Zhihua\_error.py}
  & VAFW cylinder
  & linear buckling \\
\texttt{test\_cts\_shares\_nlprebuck\_algorithm.py}
  & ---
  & source parity \\
\bottomrule
\end{tabular}
\end{center}

\noindent
The last of these compares the CTS and \texttt{newton\_raphson} modules
character by character over the regions that implement the algorithm of this
document.  The algorithm was spliced from one into the other rather than
retyped, and that test is what keeps the two from diverging silently; the
mesh generation, the per-integration-point $ABD$ of a variable-stiffness
laminate, and the CTS-specific commentary are deliberately excluded from the
comparison.

%==============================================================================
\bibliographystyle{plainnat}
\bibliography{nlprebuck_implementation}

\end{document}
